% --- Archon physics formalization source begin ---
% archon:physics
% archon:covers PhyXMiniProblems/problem_phyx_mini_0836.lean
% archon:source-report reports/phyx_mini/problem_phyx_mini_0836.source.json

\chapter{Physics problem phyx\_mini\_0836}
\label{ch:PhyXMiniProblems_problem_phyx_mini_0836}

\paragraph{Problem source.}
\#\# Physical scenario

An electron is launched at a \textbackslash{}( 45\textasciicircum{}\textbackslash{}circ \textbackslash{}) angle and a speed of \textbackslash{}( 5.0 	imes 10\textasciicircum{}6\textbackslash{},\textbackslash{}mathrm\{m/s\} \textbackslash{}) from the positive plate of the parallel-plate capacitor shown in figure. The electron lands \textbackslash{}( 4.0\textbackslash{},\textbackslash{}mathrm\{cm\} \textbackslash{}) away.

\#\# Question

What is the electric field strength inside the capacitor?

\#\# Answer choices

- A: A: 1.3 \textbackslash{}times 10\textasciicircum{}\{5\}N/C

- B: B: 5.3 \textbackslash{}times 10\textasciicircum{}\{4\}N/C

- C: C: 1.0 \textbackslash{}times 10\textasciicircum{}\{5\}N/C

- D: D: 3.6 \textbackslash{}times 10\textasciicircum{}\{3\}N/C

\#\# Figure caption (auxiliary; use the image as primary evidence)

The image shows a physics diagram with two parallel plates. The top plate has no markings, while the bottom plate is depicted with positive charges along its surface. A projectile motion is illustrated between the plates.

Key elements include:

1. **Plates**: 
   - Two horizontal plates are shown, one above the other. 
   - The bottom plate has red plus signs, indicating positive charges.

2. **Projectile Motion**:
   - A green arrow, labeled \textbackslash{}( \textbackslash{}vec\{v\}\_0 \textbackslash{}), represents the initial velocity of a projectile.
   - The arrow is directed upward at a 45-degree angle from the horizontal.
   - A dashed blue curve represents the projectile's trajectory.

3. **Angle**:
   - A 45-degree angle is marked between the initial velocity vector and the horizontal.

4. **Distance**:
   - A horizontal line with "4.0 cm" below it indicates the horizontal distance covered by the projectile.

The diagram likely illustrates a scenario of projectile motion in the presence of an electric field, influenced by the charged plates.

\#\# Dataset metadata

Electrostatics; Physical Model Grounding Reasoning, Multi-Formula Reasoning

\paragraph{Recorded answer/context.}
D: D: 3.6 \textbackslash{}times 10\textasciicircum{}\{3\}N/C

\paragraph{Figure/image path.}
/root/proposal\_for\_physic/science-mango/phyx\_mini\_run/phyx\_data/test\_image/836.png

\paragraph{Formalization target.}
create a compiling Lean file with sorry bodies at `PhyXMiniProblems/problem\_phyx\_mini\_0836.lean`. The Lean declarations must preserve the physical quantities, dimensions or dimensional roles, figure labels, governing-law hypotheses, and final relation expressed by this problem.
Use Mathlib/Physlib names found through LeanExplore where available. If a physics API is missing, introduce faithful local abstractions rather than scalar placeholder aliases.

\begin{theorem}[Physics formalization target]
\label{thm:physics:phyx_mini_0836:target}
\lean{PhyXMiniProblems.ProblemPhyXMini0836.problem_phyx_mini_0836}
\uses{def:physics:phyx-mini-0836:phyxminiproblems-problemphyxmini0836-electricfieldstrengthinnewtonspercoulomb, def:physics:phyx-mini-0836:phyxminiproblems-problemphyxmini0836-capacitorelectronsetup, def:physics:phyx-mini-0836:phyxminiproblems-problemphyxmini0836-matchesproblemandfigurereadouts, def:physics:phyx-mini-0836:phyxminiproblems-problemphyxmini0836-useselectronreferenceconstants, def:physics:phyx-mini-0836:phyxminiproblems-problemphyxmini0836-hasphysicalparameters, def:physics:phyx-mini-0836:phyxminiproblems-problemphyxmini0836-satisfiesuniformparallelplatefieldmodel, def:physics:phyx-mini-0836:phyxminiproblems-problemphyxmini0836-satisfiesconstantelectricaccelerationkinematics, def:physics:phyx-mini-0836:phyxminiproblems-problemphyxmini0836-satisfieselectricforcelaw, def:physics:phyx-mini-0836:phyxminiproblems-problemphyxmini0836-recordeddatasetanswer, def:physics:phyx-mini-0836:phyxminiproblems-problemphyxmini0836-answermatcheselectricfield}
The assigned autoformalize agent should translate this physics problem into Lean declarations in the covered file, with theorem and lemma proof bodies written as `by sorry`.
\end{theorem}
\begin{proof}
This is an autoformalization task, not a proof task. Produce statements that can later be proved by the physics prover without weakening the physical model.
\end{proof}

% --- Archon declaration topology begin ---
\section{Declaration topology}
\begin{definition}[acceleration Dimension]
  \label{def:physics:phyx-mini-0836:phyxminiproblems-problemphyxmini0836-accelerationdimension}
  \lean{PhyXMiniProblems.ProblemPhyXMini0836.accelerationDimension}
  Acceleration has dimension length divided by time squared
\end{definition}
\begin{proof}
  This is the declared model definition of acceleration Dimension.
\end{proof}
\begin{definition}[electric Field Strength Dimension]
  \label{def:physics:phyx-mini-0836:phyxminiproblems-problemphyxmini0836-electricfieldstrengthdimension}
  \lean{PhyXMiniProblems.ProblemPhyXMini0836.electricFieldStrengthDimension}
  Electric-field strength has dimension force divided by charge
\end{definition}
\begin{proof}
  This is the declared model definition of electric Field Strength Dimension.
\end{proof}
\begin{definition}[Length Quantity]
  \label{def:physics:phyx-mini-0836:phyxminiproblems-problemphyxmini0836-lengthquantity}
  \lean{PhyXMiniProblems.ProblemPhyXMini0836.LengthQuantity}
  A nonnegative, unit-independent physical length
\end{definition}
\begin{proof}
  This is the declared model definition of Length Quantity.
\end{proof}
\begin{definition}[Time Quantity]
  \label{def:physics:phyx-mini-0836:phyxminiproblems-problemphyxmini0836-timequantity}
  \lean{PhyXMiniProblems.ProblemPhyXMini0836.TimeQuantity}
  A nonnegative, unit-independent physical duration
\end{definition}
\begin{proof}
  This is the declared model definition of Time Quantity.
\end{proof}
\begin{definition}[Speed Quantity]
  \label{def:physics:phyx-mini-0836:phyxminiproblems-problemphyxmini0836-speedquantity}
  \lean{PhyXMiniProblems.ProblemPhyXMini0836.SpeedQuantity}
  A nonnegative, unit-independent physical speed
\end{definition}
\begin{proof}
  This is the declared model definition of Speed Quantity.
\end{proof}
\begin{definition}[Mass Quantity]
  \label{def:physics:phyx-mini-0836:phyxminiproblems-problemphyxmini0836-massquantity}
  \lean{PhyXMiniProblems.ProblemPhyXMini0836.MassQuantity}
  A nonnegative, unit-independent physical mass
\end{definition}
\begin{proof}
  This is the declared model definition of Mass Quantity.
\end{proof}
\begin{definition}[Charge Magnitude Quantity]
  \label{def:physics:phyx-mini-0836:phyxminiproblems-problemphyxmini0836-chargemagnitudequantity}
  \lean{PhyXMiniProblems.ProblemPhyXMini0836.ChargeMagnitudeQuantity}
  A nonnegative, unit-independent electric-charge magnitude
\end{definition}
\begin{proof}
  This is the declared model definition of Charge Magnitude Quantity.
\end{proof}
\begin{definition}[Acceleration Quantity]
  \label{def:physics:phyx-mini-0836:phyxminiproblems-problemphyxmini0836-accelerationquantity}
  \lean{PhyXMiniProblems.ProblemPhyXMini0836.AccelerationQuantity}
  \uses{def:physics:phyx-mini-0836:phyxminiproblems-problemphyxmini0836-accelerationdimension}
  A nonnegative, unit-independent acceleration magnitude
\end{definition}
\begin{proof}
  This is the declared model definition of Acceleration Quantity.
\end{proof}
\begin{definition}[Electric Field Strength Quantity]
  \label{def:physics:phyx-mini-0836:phyxminiproblems-problemphyxmini0836-electricfieldstrengthquantity}
  \lean{PhyXMiniProblems.ProblemPhyXMini0836.ElectricFieldStrengthQuantity}
  \uses{def:physics:phyx-mini-0836:phyxminiproblems-problemphyxmini0836-electricfieldstrengthdimension}
  A nonnegative, unit-independent electric-field magnitude
\end{definition}
\begin{proof}
  This is the declared model definition of Electric Field Strength Quantity.
\end{proof}
\begin{definition}[length Readout]
  \label{def:physics:phyx-mini-0836:phyxminiproblems-problemphyxmini0836-lengthreadout}
  \lean{PhyXMiniProblems.ProblemPhyXMini0836.lengthReadout}
  \uses{def:physics:phyx-mini-0836:phyxminiproblems-problemphyxmini0836-lengthquantity}
  Read a physical length in a selected named unit
\end{definition}
\begin{proof}
  This is the declared model definition of length Readout.
\end{proof}
\begin{definition}[time Readout]
  \label{def:physics:phyx-mini-0836:phyxminiproblems-problemphyxmini0836-timereadout}
  \lean{PhyXMiniProblems.ProblemPhyXMini0836.timeReadout}
  \uses{def:physics:phyx-mini-0836:phyxminiproblems-problemphyxmini0836-timequantity}
  Read a physical duration in a selected named unit
\end{definition}
\begin{proof}
  This is the declared model definition of time Readout.
\end{proof}
\begin{definition}[speed Readout]
  \label{def:physics:phyx-mini-0836:phyxminiproblems-problemphyxmini0836-speedreadout}
  \lean{PhyXMiniProblems.ProblemPhyXMini0836.speedReadout}
  \uses{def:physics:phyx-mini-0836:phyxminiproblems-problemphyxmini0836-speedquantity}
  Read a speed in coherent selected length and time units
\end{definition}
\begin{proof}
  This is the declared model definition of speed Readout.
\end{proof}
\begin{definition}[mass Readout]
  \label{def:physics:phyx-mini-0836:phyxminiproblems-problemphyxmini0836-massreadout}
  \lean{PhyXMiniProblems.ProblemPhyXMini0836.massReadout}
  \uses{def:physics:phyx-mini-0836:phyxminiproblems-problemphyxmini0836-massquantity}
  Read a physical mass in a selected named unit
\end{definition}
\begin{proof}
  This is the declared model definition of mass Readout.
\end{proof}
\begin{definition}[charge Readout]
  \label{def:physics:phyx-mini-0836:phyxminiproblems-problemphyxmini0836-chargereadout}
  \lean{PhyXMiniProblems.ProblemPhyXMini0836.chargeReadout}
  \uses{def:physics:phyx-mini-0836:phyxminiproblems-problemphyxmini0836-chargemagnitudequantity}
  Read an electric-charge magnitude in a selected named unit
\end{definition}
\begin{proof}
  This is the declared model definition of charge Readout.
\end{proof}
\begin{definition}[acceleration Readout]
  \label{def:physics:phyx-mini-0836:phyxminiproblems-problemphyxmini0836-accelerationreadout}
  \lean{PhyXMiniProblems.ProblemPhyXMini0836.accelerationReadout}
  \uses{def:physics:phyx-mini-0836:phyxminiproblems-problemphyxmini0836-accelerationquantity}
  Read an acceleration in coherent selected length and time units
\end{definition}
\begin{proof}
  This is the declared model definition of acceleration Readout.
\end{proof}
\begin{definition}[electric Field Strength In Newtons Per Coulomb]
  \label{def:physics:phyx-mini-0836:phyxminiproblems-problemphyxmini0836-electricfieldstrengthinnewtonspercoulomb}
  \lean{PhyXMiniProblems.ProblemPhyXMini0836.electricFieldStrengthInNewtonsPerCoulomb}
  \uses{def:physics:phyx-mini-0836:phyxminiproblems-problemphyxmini0836-electricfieldstrengthquantity}
  Coherent-SI readout of an electric-field strength, in newtons/coulomb
\end{definition}
\begin{proof}
  This is the declared model definition of electric Field Strength In Newtons Per Coulomb.
\end{proof}
\begin{definition}[length In Meters]
  \label{def:physics:phyx-mini-0836:phyxminiproblems-problemphyxmini0836-lengthinmeters}
  \lean{PhyXMiniProblems.ProblemPhyXMini0836.lengthInMeters}
  \uses{def:physics:phyx-mini-0836:phyxminiproblems-problemphyxmini0836-lengthquantity, def:physics:phyx-mini-0836:phyxminiproblems-problemphyxmini0836-lengthreadout}
  Metre readout of a physical length
\end{definition}
\begin{proof}
  This is the declared model definition of length In Meters.
\end{proof}
\begin{definition}[length In Centimeters]
  \label{def:physics:phyx-mini-0836:phyxminiproblems-problemphyxmini0836-lengthincentimeters}
  \lean{PhyXMiniProblems.ProblemPhyXMini0836.lengthInCentimeters}
  \uses{def:physics:phyx-mini-0836:phyxminiproblems-problemphyxmini0836-lengthquantity, def:physics:phyx-mini-0836:phyxminiproblems-problemphyxmini0836-lengthreadout}
  Centimetre readout used by the figure's 4.0 cm label
\end{definition}
\begin{proof}
  This is the declared model definition of length In Centimeters.
\end{proof}
\begin{definition}[time In Seconds]
  \label{def:physics:phyx-mini-0836:phyxminiproblems-problemphyxmini0836-timeinseconds}
  \lean{PhyXMiniProblems.ProblemPhyXMini0836.timeInSeconds}
  \uses{def:physics:phyx-mini-0836:phyxminiproblems-problemphyxmini0836-timequantity, def:physics:phyx-mini-0836:phyxminiproblems-problemphyxmini0836-timereadout}
  Second readout of a physical duration
\end{definition}
\begin{proof}
  This is the declared model definition of time In Seconds.
\end{proof}
\begin{definition}[speed In Meters Per Second]
  \label{def:physics:phyx-mini-0836:phyxminiproblems-problemphyxmini0836-speedinmeterspersecond}
  \lean{PhyXMiniProblems.ProblemPhyXMini0836.speedInMetersPerSecond}
  \uses{def:physics:phyx-mini-0836:phyxminiproblems-problemphyxmini0836-speedquantity, def:physics:phyx-mini-0836:phyxminiproblems-problemphyxmini0836-speedreadout}
  Metre-per-second readout of a physical speed
\end{definition}
\begin{proof}
  This is the declared model definition of speed In Meters Per Second.
\end{proof}
\begin{definition}[mass In Kilograms]
  \label{def:physics:phyx-mini-0836:phyxminiproblems-problemphyxmini0836-massinkilograms}
  \lean{PhyXMiniProblems.ProblemPhyXMini0836.massInKilograms}
  \uses{def:physics:phyx-mini-0836:phyxminiproblems-problemphyxmini0836-massquantity, def:physics:phyx-mini-0836:phyxminiproblems-problemphyxmini0836-massreadout}
  Kilogram readout of a physical mass
\end{definition}
\begin{proof}
  This is the declared model definition of mass In Kilograms.
\end{proof}
\begin{definition}[charge In Coulombs]
  \label{def:physics:phyx-mini-0836:phyxminiproblems-problemphyxmini0836-chargeincoulombs}
  \lean{PhyXMiniProblems.ProblemPhyXMini0836.chargeInCoulombs}
  \uses{def:physics:phyx-mini-0836:phyxminiproblems-problemphyxmini0836-chargemagnitudequantity, def:physics:phyx-mini-0836:phyxminiproblems-problemphyxmini0836-chargereadout}
  Coulomb readout of an electric-charge magnitude
\end{definition}
\begin{proof}
  This is the declared model definition of charge In Coulombs.
\end{proof}
\begin{definition}[acceleration In Meters Per Second Squared]
  \label{def:physics:phyx-mini-0836:phyxminiproblems-problemphyxmini0836-accelerationinmeterspersecondsquared}
  \lean{PhyXMiniProblems.ProblemPhyXMini0836.accelerationInMetersPerSecondSquared}
  \uses{def:physics:phyx-mini-0836:phyxminiproblems-problemphyxmini0836-accelerationquantity, def:physics:phyx-mini-0836:phyxminiproblems-problemphyxmini0836-accelerationreadout}
  Metre-per-second-squared readout of an acceleration magnitude
\end{definition}
\begin{proof}
  This is the declared model definition of acceleration In Meters Per Second Squared.
\end{proof}
\begin{definition}[elementary Charge In Coulombs]
  \label{def:physics:phyx-mini-0836:phyxminiproblems-problemphyxmini0836-elementarychargeincoulombs}
  \lean{PhyXMiniProblems.ProblemPhyXMini0836.elementaryChargeInCoulombs}
  Physlib's elementary-charge unit expressed as a scalar number of coulombs. The electron charge in the setup remains a dimensionful physical quantity
\end{definition}
\begin{proof}
  This is the declared model definition of elementary Charge In Coulombs.
\end{proof}
\begin{definition}[electron Rest Mass In Kilograms]
  \label{def:physics:phyx-mini-0836:phyxminiproblems-problemphyxmini0836-electronrestmassinkilograms}
  \lean{PhyXMiniProblems.ProblemPhyXMini0836.electronRestMassInKilograms}
  CODATA electron rest-mass readout used by the model, in kilograms
\end{definition}
\begin{proof}
  This is the declared model definition of electron Rest Mass In Kilograms.
\end{proof}
\begin{definition}[Particle Species]
  \label{def:physics:phyx-mini-0836:phyxminiproblems-problemphyxmini0836-particlespecies}
  \lean{PhyXMiniProblems.ProblemPhyXMini0836.ParticleSpecies}
  Particle species distinguished by the problem statement
\end{definition}
\begin{proof}
  This is the declared model definition of Particle Species.
\end{proof}
\begin{definition}[Charge Sign]
  \label{def:physics:phyx-mini-0836:phyxminiproblems-problemphyxmini0836-chargesign}
  \lean{PhyXMiniProblems.ProblemPhyXMini0836.ChargeSign}
  Sign of a particle's electric charge
\end{definition}
\begin{proof}
  This is the declared model definition of Charge Sign.
\end{proof}
\begin{definition}[Plate Label]
  \label{def:physics:phyx-mini-0836:phyxminiproblems-problemphyxmini0836-platelabel}
  \lean{PhyXMiniProblems.ProblemPhyXMini0836.PlateLabel}
  The two horizontal capacitor plates in the raster figure
\end{definition}
\begin{proof}
  This is the declared model definition of Plate Label.
\end{proof}
\begin{definition}[Plate Polarity]
  \label{def:physics:phyx-mini-0836:phyxminiproblems-problemphyxmini0836-platepolarity}
  \lean{PhyXMiniProblems.ProblemPhyXMini0836.PlatePolarity}
  Electrical polarity shown by the marks on a plate
\end{definition}
\begin{proof}
  This is the declared model definition of Plate Polarity.
\end{proof}
\begin{definition}[Plate Arrangement]
  \label{def:physics:phyx-mini-0836:phyxminiproblems-problemphyxmini0836-platearrangement}
  \lean{PhyXMiniProblems.ProblemPhyXMini0836.PlateArrangement}
  Relative geometry of the two capacitor plates
\end{definition}
\begin{proof}
  This is the declared model definition of Plate Arrangement.
\end{proof}
\begin{definition}[Spatial Direction]
  \label{def:physics:phyx-mini-0836:phyxminiproblems-problemphyxmini0836-spatialdirection}
  \lean{PhyXMiniProblems.ProblemPhyXMini0836.SpatialDirection}
  Directions needed to interpret the side-view diagram
\end{definition}
\begin{proof}
  This is the declared model definition of Spatial Direction.
\end{proof}
\begin{definition}[Trajectory Style]
  \label{def:physics:phyx-mini-0836:phyxminiproblems-problemphyxmini0836-trajectorystyle}
  \lean{PhyXMiniProblems.ProblemPhyXMini0836.TrajectoryStyle}
  Rendering style of the pictured electron trajectory
\end{definition}
\begin{proof}
  This is the declared model definition of Trajectory Style.
\end{proof}
\begin{definition}[Electron Motion Model]
  \label{def:physics:phyx-mini-0836:phyxminiproblems-problemphyxmini0836-electronmotionmodel}
  \lean{PhyXMiniProblems.ProblemPhyXMini0836.ElectronMotionModel}
  Idealization used while the electron is between the plates
\end{definition}
\begin{proof}
  This is the declared model definition of Electron Motion Model.
\end{proof}
\begin{definition}[Parallel Plate Figure]
  \label{def:physics:phyx-mini-0836:phyxminiproblems-problemphyxmini0836-parallelplatefigure}
  \lean{PhyXMiniProblems.ProblemPhyXMini0836.ParallelPlateFigure}
  \uses{def:physics:phyx-mini-0836:phyxminiproblems-problemphyxmini0836-lengthquantity, def:physics:phyx-mini-0836:phyxminiproblems-problemphyxmini0836-platelabel, def:physics:phyx-mini-0836:phyxminiproblems-problemphyxmini0836-platepolarity, def:physics:phyx-mini-0836:phyxminiproblems-problemphyxmini0836-platearrangement, def:physics:phyx-mini-0836:phyxminiproblems-problemphyxmini0836-spatialdirection, def:physics:phyx-mini-0836:phyxminiproblems-problemphyxmini0836-trajectorystyle}
  Typed content read from the supplied raster. The horizontal range label is a physical length; its numerical centimetre readout is stated separately
\end{definition}
\begin{proof}
  This is the declared model definition of Parallel Plate Figure.
\end{proof}
\begin{definition}[Capacitor Electron Setup]
  \label{def:physics:phyx-mini-0836:phyxminiproblems-problemphyxmini0836-capacitorelectronsetup}
  \lean{PhyXMiniProblems.ProblemPhyXMini0836.CapacitorElectronSetup}
  \uses{def:physics:phyx-mini-0836:phyxminiproblems-problemphyxmini0836-lengthquantity, def:physics:phyx-mini-0836:phyxminiproblems-problemphyxmini0836-timequantity, def:physics:phyx-mini-0836:phyxminiproblems-problemphyxmini0836-speedquantity, def:physics:phyx-mini-0836:phyxminiproblems-problemphyxmini0836-massquantity, def:physics:phyx-mini-0836:phyxminiproblems-problemphyxmini0836-chargemagnitudequantity, def:physics:phyx-mini-0836:phyxminiproblems-problemphyxmini0836-accelerationquantity, def:physics:phyx-mini-0836:phyxminiproblems-problemphyxmini0836-electricfieldstrengthquantity, def:physics:phyx-mini-0836:phyxminiproblems-problemphyxmini0836-particlespecies, def:physics:phyx-mini-0836:phyxminiproblems-problemphyxmini0836-chargesign, def:physics:phyx-mini-0836:phyxminiproblems-problemphyxmini0836-spatialdirection, def:physics:phyx-mini-0836:phyxminiproblems-problemphyxmini0836-electronmotionmodel, def:physics:phyx-mini-0836:phyxminiproblems-problemphyxmini0836-parallelplatefigure}
  The independent physical objects and observables. In particular, neither electricFieldStrength nor accelerationMagnitude is defined from the recorded answer or from the desired 3.6 * 10\textasciicircum{}3 N/C value
\end{definition}
\begin{proof}
  This is the declared model definition of Capacitor Electron Setup.
\end{proof}
\begin{definition}[Matches Problem And Figure Readouts]
  \label{def:physics:phyx-mini-0836:phyxminiproblems-problemphyxmini0836-matchesproblemandfigurereadouts}
  \lean{PhyXMiniProblems.ProblemPhyXMini0836.MatchesProblemAndFigureReadouts}
  \uses{def:physics:phyx-mini-0836:phyxminiproblems-problemphyxmini0836-lengthinmeters, def:physics:phyx-mini-0836:phyxminiproblems-problemphyxmini0836-lengthincentimeters, def:physics:phyx-mini-0836:phyxminiproblems-problemphyxmini0836-speedinmeterspersecond, def:physics:phyx-mini-0836:phyxminiproblems-problemphyxmini0836-capacitorelectronsetup}
  Problem-statement and primary-image evidence. The raster shows negative marks on the upper plate and positive marks on the lower plate, even though the auxiliary caption says that the upper plate is unmarked. No field-strength answer occurs in these data
\end{definition}
\begin{proof}
  This is the declared model definition of Matches Problem And Figure Readouts.
\end{proof}
\begin{definition}[Uses Electron Reference Constants]
  \label{def:physics:phyx-mini-0836:phyxminiproblems-problemphyxmini0836-useselectronreferenceconstants}
  \lean{PhyXMiniProblems.ProblemPhyXMini0836.UsesElectronReferenceConstants}
  \uses{def:physics:phyx-mini-0836:phyxminiproblems-problemphyxmini0836-massinkilograms, def:physics:phyx-mini-0836:phyxminiproblems-problemphyxmini0836-chargeincoulombs, def:physics:phyx-mini-0836:phyxminiproblems-problemphyxmini0836-elementarychargeincoulombs, def:physics:phyx-mini-0836:phyxminiproblems-problemphyxmini0836-electronrestmassinkilograms, def:physics:phyx-mini-0836:phyxminiproblems-problemphyxmini0836-capacitorelectronsetup}
  Standard reference calibrations needed to turn the electron's acceleration into an electric-field strength. These are independent particle data, not a statement of the requested field value
\end{definition}
\begin{proof}
  This is the declared model definition of Uses Electron Reference Constants.
\end{proof}
\begin{definition}[Has Physical Parameters]
  \label{def:physics:phyx-mini-0836:phyxminiproblems-problemphyxmini0836-hasphysicalparameters}
  \lean{PhyXMiniProblems.ProblemPhyXMini0836.HasPhysicalParameters}
  \uses{def:physics:phyx-mini-0836:phyxminiproblems-problemphyxmini0836-electricfieldstrengthinnewtonspercoulomb, def:physics:phyx-mini-0836:phyxminiproblems-problemphyxmini0836-lengthinmeters, def:physics:phyx-mini-0836:phyxminiproblems-problemphyxmini0836-timeinseconds, def:physics:phyx-mini-0836:phyxminiproblems-problemphyxmini0836-speedinmeterspersecond, def:physics:phyx-mini-0836:phyxminiproblems-problemphyxmini0836-massinkilograms, def:physics:phyx-mini-0836:phyxminiproblems-problemphyxmini0836-chargeincoulombs, def:physics:phyx-mini-0836:phyxminiproblems-problemphyxmini0836-accelerationinmeterspersecondsquared, def:physics:phyx-mini-0836:phyxminiproblems-problemphyxmini0836-capacitorelectronsetup}
  Positivity and non-vacuity conditions for the physical experiment
\end{definition}
\begin{proof}
  This is the declared model definition of Has Physical Parameters.
\end{proof}
\begin{definition}[Satisfies Uniform Parallel Plate Field Model]
  \label{def:physics:phyx-mini-0836:phyxminiproblems-problemphyxmini0836-satisfiesuniformparallelplatefieldmodel}
  \lean{PhyXMiniProblems.ProblemPhyXMini0836.SatisfiesUniformParallelPlateFieldModel}
  \uses{def:physics:phyx-mini-0836:phyxminiproblems-problemphyxmini0836-electricfieldstrengthinnewtonspercoulomb, def:physics:phyx-mini-0836:phyxminiproblems-problemphyxmini0836-capacitorelectronsetup}
  Uniform parallel-plate electrostatics: the field points from the positive lower plate toward the negative upper plate, and its scalar physical strength is the norm of Physlib's spacetime-dependent vector field throughout the gap. This law contains no numerical value for the field strength
\end{definition}
\begin{proof}
  This is the declared model definition of Satisfies Uniform Parallel Plate Field Model.
\end{proof}
\begin{definition}[Satisfies Constant Electric Acceleration Kinematics]
  \label{def:physics:phyx-mini-0836:phyxminiproblems-problemphyxmini0836-satisfiesconstantelectricaccelerationkinematics}
  \lean{PhyXMiniProblems.ProblemPhyXMini0836.SatisfiesConstantElectricAccelerationKinematics}
  \uses{def:physics:phyx-mini-0836:phyxminiproblems-problemphyxmini0836-lengthinmeters, def:physics:phyx-mini-0836:phyxminiproblems-problemphyxmini0836-timeinseconds, def:physics:phyx-mini-0836:phyxminiproblems-problemphyxmini0836-speedinmeterspersecond, def:physics:phyx-mini-0836:phyxminiproblems-problemphyxmini0836-accelerationinmeterspersecondsquared, def:physics:phyx-mini-0836:phyxminiproblems-problemphyxmini0836-capacitorelectronsetup}
  Equal-height constant-acceleration endpoint equations. Horizontally the electron moves uniformly; vertically its initial upward component is exactly cancelled by the downward electric acceleration at landing. These are the general kinematic laws, not the requested electric-field result
\end{definition}
\begin{proof}
  This is the declared model definition of Satisfies Constant Electric Acceleration Kinematics.
\end{proof}
\begin{definition}[Satisfies Electric Force Law]
  \label{def:physics:phyx-mini-0836:phyxminiproblems-problemphyxmini0836-satisfieselectricforcelaw}
  \lean{PhyXMiniProblems.ProblemPhyXMini0836.SatisfiesElectricForceLaw}
  \uses{def:physics:phyx-mini-0836:phyxminiproblems-problemphyxmini0836-electricfieldstrengthinnewtonspercoulomb, def:physics:phyx-mini-0836:phyxminiproblems-problemphyxmini0836-massinkilograms, def:physics:phyx-mini-0836:phyxminiproblems-problemphyxmini0836-chargeincoulombs, def:physics:phyx-mini-0836:phyxminiproblems-problemphyxmini0836-accelerationinmeterspersecondsquared, def:physics:phyx-mini-0836:phyxminiproblems-problemphyxmini0836-capacitorelectronsetup}
  The magnitude form of the electric part of the Lorentz-force law together with Newton's second law: m a = |q| E. Directional information for the negative electron is kept in the parallel-plate field model above
\end{definition}
\begin{proof}
  This is the declared model definition of Satisfies Electric Force Law.
\end{proof}
\begin{lemma}[electric Field Strength eq mass mul speed sq div charge mul range]
  \label{lem:physics:phyx-mini-0836:phyxminiproblems-problemphyxmini0836-electricfieldstrength-eq-mass-mul-speed-sq-div-charge-mul-range}
  \lean{PhyXMiniProblems.ProblemPhyXMini0836.electricFieldStrength_eq_mass_mul_speed_sq_div_charge_mul_range}
  \uses{def:physics:phyx-mini-0836:phyxminiproblems-problemphyxmini0836-electricfieldstrengthinnewtonspercoulomb, def:physics:phyx-mini-0836:phyxminiproblems-problemphyxmini0836-lengthinmeters, def:physics:phyx-mini-0836:phyxminiproblems-problemphyxmini0836-speedinmeterspersecond, def:physics:phyx-mini-0836:phyxminiproblems-problemphyxmini0836-massinkilograms, def:physics:phyx-mini-0836:phyxminiproblems-problemphyxmini0836-chargeincoulombs, def:physics:phyx-mini-0836:phyxminiproblems-problemphyxmini0836-capacitorelectronsetup, def:physics:phyx-mini-0836:phyxminiproblems-problemphyxmini0836-matchesproblemandfigurereadouts, def:physics:phyx-mini-0836:phyxminiproblems-problemphyxmini0836-hasphysicalparameters, def:physics:phyx-mini-0836:phyxminiproblems-problemphyxmini0836-satisfiesconstantelectricaccelerationkinematics, def:physics:phyx-mini-0836:phyxminiproblems-problemphyxmini0836-satisfieselectricforcelaw}
  At 45°, the equal-height endpoint equations give a = v₀² / R; combining this with m a = |q| E gives the exact SI-readout relation below. The field strength remains an independent physical observable in the setup
\end{lemma}
\begin{proof}
  The conclusion follows from the governing relations and figure readouts named in the statement of electric Field Strength eq mass mul speed sq div charge mul range.
\end{proof}
\begin{definition}[Answer Choice]
  \label{def:physics:phyx-mini-0836:phyxminiproblems-problemphyxmini0836-answerchoice}
  \lean{PhyXMiniProblems.ProblemPhyXMini0836.AnswerChoice}
  Labels of the four field-strength choices printed in the question
\end{definition}
\begin{proof}
  This is the declared model definition of Answer Choice.
\end{proof}
\begin{definition}[answer Strength In Newtons Per Coulomb]
  \label{def:physics:phyx-mini-0836:phyxminiproblems-problemphyxmini0836-answerstrengthinnewtonspercoulomb}
  \lean{PhyXMiniProblems.ProblemPhyXMini0836.answerStrengthInNewtonsPerCoulomb}
  \uses{def:physics:phyx-mini-0836:phyxminiproblems-problemphyxmini0836-answerchoice}
  Electric-field strength in newtons/coulomb printed beside each choice
\end{definition}
\begin{proof}
  This is the declared model definition of answer Strength In Newtons Per Coulomb.
\end{proof}
\begin{definition}[answer Rounding Tolerance]
  \label{def:physics:phyx-mini-0836:phyxminiproblems-problemphyxmini0836-answerroundingtolerance}
  \lean{PhyXMiniProblems.ProblemPhyXMini0836.answerRoundingTolerance}
  \uses{def:physics:phyx-mini-0836:phyxminiproblems-problemphyxmini0836-answerchoice}
  Half of one unit in the last displayed significant digit. This gives a precise meaning to matching the rounded numerical choices
\end{definition}
\begin{proof}
  This is the declared model definition of answer Rounding Tolerance.
\end{proof}
\begin{definition}[recorded Dataset Answer]
  \label{def:physics:phyx-mini-0836:phyxminiproblems-problemphyxmini0836-recordeddatasetanswer}
  \lean{PhyXMiniProblems.ProblemPhyXMini0836.recordedDatasetAnswer}
  \uses{def:physics:phyx-mini-0836:phyxminiproblems-problemphyxmini0836-answerchoice}
  The answer label recorded by the source dataset
\end{definition}
\begin{proof}
  This is the declared model definition of recorded Dataset Answer.
\end{proof}
\begin{definition}[Answer Matches Electric Field]
  \label{def:physics:phyx-mini-0836:phyxminiproblems-problemphyxmini0836-answermatcheselectricfield}
  \lean{PhyXMiniProblems.ProblemPhyXMini0836.AnswerMatchesElectricField}
  \uses{def:physics:phyx-mini-0836:phyxminiproblems-problemphyxmini0836-electricfieldstrengthinnewtonspercoulomb, def:physics:phyx-mini-0836:phyxminiproblems-problemphyxmini0836-capacitorelectronsetup, def:physics:phyx-mini-0836:phyxminiproblems-problemphyxmini0836-answerchoice, def:physics:phyx-mini-0836:phyxminiproblems-problemphyxmini0836-answerstrengthinnewtonspercoulomb, def:physics:phyx-mini-0836:phyxminiproblems-problemphyxmini0836-answerroundingtolerance}
  A choice matches the independently modeled physical field when the SI readout lies within half a unit of the choice's last displayed digit
\end{definition}
\begin{proof}
  This is the declared model definition of Answer Matches Electric Field.
\end{proof}
\begin{lemma}[electric Field Strength rounds to three Point Six Times Ten Cubed]
  \label{lem:physics:phyx-mini-0836:phyxminiproblems-problemphyxmini0836-electricfieldstrength-rounds-to-threepointsixtimestencubed}
  \lean{PhyXMiniProblems.ProblemPhyXMini0836.electricFieldStrength_rounds_to_threePointSixTimesTenCubed}
  \uses{def:physics:phyx-mini-0836:phyxminiproblems-problemphyxmini0836-electricfieldstrengthinnewtonspercoulomb, def:physics:phyx-mini-0836:phyxminiproblems-problemphyxmini0836-capacitorelectronsetup, def:physics:phyx-mini-0836:phyxminiproblems-problemphyxmini0836-matchesproblemandfigurereadouts, def:physics:phyx-mini-0836:phyxminiproblems-problemphyxmini0836-useselectronreferenceconstants, def:physics:phyx-mini-0836:phyxminiproblems-problemphyxmini0836-hasphysicalparameters, def:physics:phyx-mini-0836:phyxminiproblems-problemphyxmini0836-satisfiesconstantelectricaccelerationkinematics, def:physics:phyx-mini-0836:phyxminiproblems-problemphyxmini0836-satisfieselectricforcelaw}
  The launch data, CODATA electron constants, ideal endpoint kinematics, and electric-force law put the field within 50 N/C of 3.6 * 10\textasciicircum{}3 N/C, so the two-significant-digit displayed answer is choice D
\end{lemma}
\begin{proof}
  The conclusion follows from the governing relations and figure readouts named in the statement of electric Field Strength rounds to three Point Six Times Ten Cubed.
\end{proof}
% --- Archon declaration topology end ---
% --- Archon physics formalization source end ---
